\documentclass[a4paper, serif, 10pt]{moderncv}

%% moderncv
% style and color
\moderncvstyle{banking}
\moderncvcolor{black}

%% page layout/margins
\usepackage[top=2.5cm, bottom=2.5cm, left=1.5cm, right=1.5cm]{geometry}

\nopagenumbers{}

%% Babel config for Spanish language
% \usepackage[spanish]{babel} has some conflicting issues with moderncv
% so we have to use enable the following options to make it work
%
% ref:
% - https://tex.stackexchange.com/a/140161/36007
\usepackage[spanish,es-lcroman]{babel}

%% fontspec
\usepackage{fontspec}

\IfFontExistsTF{Linux Libertine}{
  \setmainfont[Ligatures={TeX, Common}, Numbers=OldStyle]{Linux Libertine}
}{}
\IfFontExistsTF{Linux Libertine O}{
  \setmainfont[Ligatures={TeX, Common}, Numbers=OldStyle]{Linux Libertine O}
}{}

%% CTeX
% CJK support, used to show CJK characters in the resume
%
% - fontset=none: disable builtin fontset but instead set the CJK font manually
% - heading=false: disable ctex heading
% - punct=kaiming: use kaiming punctuations styles for CJK
% - scheme=plain: use plain scheme, do not override \normalsize font size
% - space=auto: space settings for CJK characters
%
% ref:
% - http://ctan.mirrorcatalogs.com/language/chinese/ctex/ctex.pdf
\usepackage[UTF8, heading=false, punct=kaiming, scheme=plain, space=auto]{ctex}

\IfFontExistsTF{Noto Serif CJK SC}{
  \setCJKmainfont{Noto Serif CJK SC}
}{}
\IfFontExistsTF{Noto Sans CJK SC}{
  \setCJKsansfont{Noto Sans CJK SC}
}{}

%% line spacing
\usepackage{setspace}
\setstretch{1.125}

%% URL styling - use normal text instead of monospace
\urlstyle{same}

% Auto-underline all links
% Defer until \begin{document} so hyperref is already loaded
\AtBeginDocument{%
  \let\oldhref\href
  \renewcommand{\href}[2]{\oldhref{#1}{\underline{#2}}}%
}

%% Patch \httplink and \httpslink to handle full URLs
% The original moderncv macros always prepend http:// or https:// to URLs.
% This causes issues when the URL already has a protocol (e.g., https://example.com)
% resulting in malformed URLs like https://https://example.com.
% This patch checks if the URL already contains "://" and uses it directly if so.
% Note: We use \str_set:Nx (x-expansion) to expand macros like \@homepage first.
\ExplSyntaxOn
\str_new:N \g_yamlresume_url_str

\RenewDocumentCommand{\httpslink}{O{}m}{
  \str_set:Nx \g_yamlresume_url_str {#2}
  \str_if_in:NnTF \g_yamlresume_url_str {://}
    { \url{#2} }
    { \url{https://#2} }
}

\RenewDocumentCommand{\httplink}{O{}m}{
  \str_set:Nx \g_yamlresume_url_str {#2}
  \str_if_in:NnTF \g_yamlresume_url_str {://}
    { \url{#2} }
    { \url{http://#2} }
}
\ExplSyntaxOff

\name{Lucía Fernández Gómez}{}
\title{Ingeniera DevOps Senior | Cloud, Kubernetes y Automatización}
\phone[mobile]{+34 612 345 678}
\email{lucia.fernandez.devops@example.com}
\homepage{https://www.luciafernandez.dev}

\address{Calle de Serrano 45, Madrid, Comunidad de Madrid, España, 28001}{}{}

\extrainfo{{\small \faGithub}\ \href{https://github.com/luciaferdevops}{@luciaferdevops} {} {} {} • {} {} {} 
{\small \faLinkedin}\ \href{https://www.linkedin.com/in/lucia-fernandez-devops}{@lucia-fernandez-devops} {} {} {} • {} {} {} 
{\small \faGitlab}\ \href{https://gitlab.com/luciaferdevops}{@luciaferdevops}}

\begin{document}

\maketitle

\section{Información básica}

\cvline{}{Ingeniera DevOps con más de 8 años de experiencia diseñando, automatizando y operando plataformas cloud críticas. Especialista en infraestructura como código, contenedores y pipelines CI/CD para equipos de alto rendimiento.
%
\begin{itemize}
\item Reduje el tiempo de despliegue en un 70 \% mediante pipelines automatizados y estrategias GitOps.
\item Diseñé arquitecturas multinube (AWS, Azure y GCP) altamente disponibles, seguras y orientadas al coste.
\item Lideré la migración de cargas de trabajo monolíticas a Kubernetes, mejorando la escalabilidad y la resiliencia.
\item Promoví cultura DevOps, observabilidad y respuesta ante incidentes con equipos de desarrollo y SRE.
\end{itemize}}
\section{Educación}

\cventry{sept 2012–jul 2016}
        {Licenciatura, Ingeniería Informática, Puntuación: 8,7/10}
        {Universidad Politécnica de Madrid}
        {\url{https://www.upm.es}}
        {}
        {\begin{itemize}
\item Formación sólida en fundamentos de computación, redes y sistemas distribuidos.
\item Proyecto fin de grado sobre automatización de despliegues en entornos virtualizados.
\item Participación activa en asociaciones tecnológicas y talleres de software libre.
\end{itemize}
\textbf{Cursos}: Sistemas Operativos, Redes de Computadores, Arquitectura de Software, Bases de Datos, Seguridad Informática, Computación en la Nube}

\cventry{sept 2016–jun 2018}
        {Maestría, Ingeniería de Sistemas y Computación, Puntuación: 9,1/10}
        {Universidad Nacional de Educación a Distancia}
        {\url{https://www.uned.es}}
        {}
        {\begin{itemize}
\item Especialización en plataformas cloud, orquestación de contenedores y automatización operativa.
\item Trabajo fin de máster sobre optimización de costes en clústeres Kubernetes.
\end{itemize}
\textbf{Cursos}: Virtualización y Contenedores, Computación Distribuida, Gestión de Infraestructuras Cloud, Metodologías Ágiles y DevOps}
\section{Experiencia laboral}

\cventry{mar 2021 hasta la fecha}
        {Ingeniera DevOps Senior}
        {Telefónica Tech}
        {\url{https://www.telefonicatech.com}}
        {}
        {\begin{itemize}
\item Diseñé y mantuve plataformas Kubernetes multirregión para más de 200 microservicios en producción.
\item Implementé infraestructura como código con Terraform y Ansible, reduciendo la deriva de configuración en un 85 \%.
\item Construí pipelines CI/CD con GitLab CI y Argo CD, logrando despliegues progresivos con canary releases.
\item Definí SLOs, alertas y cuadros de mando con Prometheus, Grafana y OpenTelemetry.
\item Automaticé la gestión de secretos con HashiCorp Vault y políticas de mínimo privilegio.
\end{itemize}
\textbf{Palabras clave}: Kubernetes, Terraform, GitOps, Prometheus, AWS}

\cventry{jul 2018–feb 2021}
        {Ingeniera de Plataforma Cloud}
        {BBVA}
        {\url{https://www.bbva.com}}
        {}
        {\begin{itemize}
\item Migré servicios críticos a AWS con arquitecturas sin servidor y basadas en contenedores.
\item Desarrollé módulos reutilizables de Terraform para equipos de desarrollo, acelerando el aprovisionamiento.
\item Implanté monitorización centralizada con ELK Stack y alertas proactivas para reducir el MTTR.
\item Participé en guardias 24x7 y en análisis post-mortem de incidentes en producción.
\end{itemize}
\textbf{Palabras clave}: AWS, Docker, Jenkins, ELK, Python}

\cventry{sept 2016–jun 2018}
        {Administradora de Sistemas}
        {Indra Sistemas}
        {\url{https://www.indracompany.com}}
        {}
        {\begin{itemize}
\item Administré entornos Linux y Windows Server en infraestructuras de alta disponibilidad.
\item Automaticé tareas de despliegue y mantenimiento con scripts Bash y Python.
\item Configuré balanceadores, firewalls y monitorización con Nagios y Zabbix.
\end{itemize}
\textbf{Palabras clave}: Linux, Bash, Nagios, Zabbix, Redes}
\section{Idiomas}

\cvline{Español}{Competencia nativa o bilingüe \hfill \textbf{Palabras clave}: Español nativo}
\cvline{Inglés}{Competencia profesional plena \hfill \textbf{Palabras clave}: C1, TOEFL 110}
\cvline{Catalán}{Competencia limitada de trabajo \hfill \textbf{Palabras clave}: Conversación básica}
\section{Competencias}

\cvline{Computación en la nube}{Experto \hfill \textbf{Palabras clave}: AWS, Azure, GCP, Terraform}
\cvline{Contenedores y orquestación}{Experto \hfill \textbf{Palabras clave}: Kubernetes, Docker, Helm, Istio}
\cvline{Automatización y CI/CD}{Avanzado \hfill \textbf{Palabras clave}: GitLab CI, Jenkins, Argo CD, Ansible}
\cvline{Observabilidad}{Avanzado \hfill \textbf{Palabras clave}: Prometheus, Grafana, OpenTelemetry, ELK}
\cvline{Seguridad y cumplimiento}{Avanzado \hfill \textbf{Palabras clave}: Vault, Trivy, OPA, IAM}
\cvline{Programación y scripting}{Avanzado \hfill \textbf{Palabras clave}: Python, Bash, Go, YAML}
\section{Premios}

\cventry{dic 2022}
        {Telefónica Tech}
        {Premio a la Excelencia Técnica}
        {}
        {}
        {Reconocimiento por liderar la plataforma Kubernetes multirregión y reducir los incidentes críticos en un 40 \% durante 2022.}

\cventry{jun 2019}
        {BBVA}
        {Mención de Honor en Innovación Cloud}
        {}
        {}
        {Distinción por el diseño de módulos Terraform reutilizables que aceleraron la migración de servicios críticos a AWS.}
\section{Certificados}

\cventry{may 2022}
        {Amazon Web Services}
        {AWS Certified DevOps Engineer – Professional}
        {\url{https://aws.amazon.com/certification/}}
        {}
        {}

\cventry{nov 2021}
        {Cloud Native Computing Foundation}
        {Certified Kubernetes Administrator (CKA)}
        {\url{https://www.cncf.io/certification/cka/}}
        {}
        {}

\cventry{mar 2023}
        {HashiCorp}
        {HashiCorp Certified Terraform Associate}
        {\url{https://www.hashicorp.com/certification/terraform-associate}}
        {}
        {}
\section{Publicaciones}

\cventry{oct 2023}
        {Automatización GitOps para plataformas multinube}
        {Revista Cloud \& DevOps}
        {\url{https://www.clouddevops.example.com/articulos/gitops-multinube}}
        {}
        {\begin{itemize}
\item Artículo técnico sobre patrones GitOps para desplegar aplicaciones en AWS, Azure y GCP.
\item Incluye ejemplos con Argo CD, Terraform y políticas de promoción entre entornos.
\item Analiza métricas de adopción, seguridad y reducción del tiempo de despliegue.
\end{itemize}}
\section{Proyectos}

\cventry{ene 2022–dic 2023}
        {Plataforma interna para desplegar microservicios en varios proveedores cloud mediante GitOps.}
        {Plataforma GitOps Multinube}
        {\url{https://github.com/luciaferdevops/gitops-multinube}}
        {}
        {\begin{itemize}
\item Implementé flujos de promoción automática entre desarrollo, preproducción y producción.
\item Integré validaciones de seguridad y coste en los pipelines antes del despliegue.
\item Reduje el tiempo medio de despliegue de 45 a 12 minutos por servicio.
\end{itemize}
\textbf{Palabras clave}: GitOps, Argo CD, Terraform, Kubernetes}

\cventry{mar 2020–nov 2021}
        {Solución de observabilidad centralizada para equipos de desarrollo y SRE.}
        {Observabilidad Unificada}
        {\url{https://github.com/luciaferdevops/observabilidad-unificada}}
        {}
        {\begin{itemize}
\item Desplegué Prometheus, Grafana y OpenTelemetry en clústeres Kubernetes de alta disponibilidad.
\item Definí SLOs y alertas accionables, reduciendo el ruido de alertas en un 60 \%.
\item Documenté runbooks y formé a equipos en investigación de incidentes.
\end{itemize}
\textbf{Palabras clave}: Prometheus, Grafana, OpenTelemetry, SRE}
\section{Intereses}

\cvline{Deportes}{Senderismo, Ciclismo, Natación}
\cvline{Música}{Guitarra, Piano}
\cvline{Comunidad tech}{Meetups, Mentoría, Open Source}
\section{Voluntariado}

\cventry{sept 2019 hasta la fecha}
        {Mentora Técnica}
        {Mujeres en DevOps Madrid}
        {\url{https://www.mujeresendevops.example.com}}
        {}
        {\begin{itemize}
\item Mentoricé a profesionales que inician carrera en DevOps y cloud.
\item Impartí talleres prácticos de Docker, Kubernetes y CI/CD.
\item Ayudé a organizar encuentros mensuales con charlas técnicas.
\end{itemize}}
    

\cventry{ene 2018–dic 2019}
        {Voluntaria en Digitalización}
        {Cáritas Madrid}
        {\url{https://www.caritasmadrid.org}}
        {}
        {\begin{itemize}
\item Automaticé la gestión de inventario y donaciones con hojas de cálculo y scripts.
\item Formé al personal en herramientas digitales básicas y buenas prácticas de seguridad.
\item Colaboré en la migración de datos a una plataforma en la nube.
\end{itemize}}
    
\end{document}